\documentclass[10pt,letterpaper]{article} % exam class defaults to 10pt on letterpaper
\usepackage[margin=1in]{geometry}         % exam class defaults to 1-inch margins
\usepackage{graphicx} 
\usepackage{fancyhdr}
\usepackage{amssymb}
\usepackage{amsmath}
\usepackage{mathtools}
\usepackage[most]{tcolorbox}
\usepackage[dvipsnames]{xcolor}
\usepackage{blindtext}
\usepackage{amsthm}
\usepackage{upgreek}
\usepackage{hyperref}
\usepackage{tikz}
\usetikzlibrary{positioning}
%making an edit, want to see if auto script commits it%
%making another edit to test the commit bot

\setlength{\parindent}{0pt}
\setlength{\parskip}{1ex plus 0.5ex minus 0.2ex} % Matches the paragraph spacing of exam
\hbadness=10000 % Place in your preamble


\newtheorem{thm}{Theorem}[section]
\newtheorem{lem}{Lemma}[section]
\newtheorem{defn}{Definition}[section]
\newtheorem{cor}{Corollary}[section]
\newtheorem{axm}{Axiom}[section]

\tcolorboxenvironment{thm}{colback=lime, colframe=black, breakable}
\tcolorboxenvironment{lem}{colback=yellow, colframe=black, breakable}
\tcolorboxenvironment{defn}{colback=pink, colframe=black, breakable}
\tcolorboxenvironment{cor}{colback=SpringGreen, colframe=black, breakable}
\tcolorboxenvironment{axm}{colback=white, colframe=black, breakable}

% --- Header/Footer Setup ---
\fancyhf{}
\renewcommand\headrulewidth{0.5pt}
\pagestyle{fancy}
\author{Omkar Tasgaonkar}
\setlength{\headsep}{0.5in}
\fancyfoot[C]{(\thepage)}
\renewcommand\footrulewidth{0.5pt}

\title{Ordinary Differential Equations \\[0.5em] \large Through the Lens of Nets}
\author{Omkar Tasgaonkar \\ \small otasgaonkar@ucla.edu}
\date{}

\fancypagestyle{firstpage}{
    \fancyhf{}
    \fancyhead[C]{\textbf{Notes in Analysis}}
    \renewcommand\headrulewidth{0.5pt}
    \fancyfoot[C]{(\thepage)}
    \renewcommand\footrulewidth{0.5pt}
}
\begin{document}
\maketitle
I am writing this document as my own personal notes regarding analysis and differential equations. I am attempting to cover nets and compactness rigorously in order to precisely explore Integration and Arzela-Ascoli (and related theorems/statements), which are key to the study of Ordinary Differential Equations.\\

Another topic I plan to treat thoroughly is the topic of total variation and bounded variation, which is recurrent with regards to the Riemann-Stieltjes integral.\\

\begin{tcolorbox}[colback=white, colframe=black, arc=0mm]
Textbooks Used to Develop these Notes:
\begin{enumerate}
\item \textit{\textbf{Principles of Mathematical Analysis} - Walter Rudin}
\item \textit{\textbf{Mathematical Analysis} - Tom Apostol} (for integration)
\item \textit{\textbf{General Topology} - John Kelley} (for nets)
\item \textit{\textbf{Theory of Ordinary Differential Equations} - Earl Coddington, Norman Levinson}
\item \textit{\textbf{Measure, Integration, and Real Analysis} - Sheldon Axler} (for measure theory)
\end{enumerate}
\end{tcolorbox}
\thispagestyle{firstpage}
\newpage
\begin{tcolorbox}[colback=white, colframe=black, arc=0mm]
\tableofcontents
\end{tcolorbox}
\textbf{NOTE:} Any sections/subsections marked with a (*) are deviations from the flow of content and have been added to this document as extensions of introduced theory.
\setcounter{section}{-1}

\pagebreak

\section{Nets and Sequences}

\subsection{First Countability}
What has been giving us sequential equivalence in $\mathbb{R}$ is the first countability of $\mathbb{R}$ (which means around any point $x$, there exists a countable basis). We used this to prove the more general properties implied sequential properties (by using the first countable property to construct/apply subsequences).\\

To state more rigorously, a space $X$ is first countable if, given some point $x \in X$ and countably infinite set of local nested neighborhoods of $x$ which we denote as $\mathcal{B}_n \coloneq \{B_n : n \in \mathbb{N}\}$, then:
\begin{equation*}
\forall U_x, \exists n_0 \in \mathbb{N}, \forall n \geqslant n_0: B_n \subseteq U_x
\end{equation*}
where $U_x$ is any arbitrary neighborhood of $x$.\\

All metric spaces are first countable; but there are some topological spaces that are not first countable. Consider the following topology on $\mathbb{R}$:
\begin{equation*}
\uptau \coloneq \{\varnothing\} \cup \{\mathbb{R} \setminus F: F \ \text{is finite}\}
\end{equation*}
Now take some point $x \in \mathbb{R}$ and assume it has some countable basis $\mathcal{B}_x \coloneq \{B_n\}_{n \in \mathbb{N}}$, in which case:
\begin{equation*}
\forall n \in \mathbb{N}: B_n = \mathbb{R}\setminus F_n
\end{equation*}
Then, the following set should be countably infinite:
\begin{equation*}
\bigcup_{n \in \mathbb{N}}(\mathbb{R}\setminus B_n)
\end{equation*}
since each $F_n$ is finite. We note that:
\begin{equation*}
\bigcup_{n \in \mathbb{N}}(\mathbb{R}\setminus B_n) = \mathbb{R} \setminus \bigcap_{n \in \mathbb{N}} B_n = \mathbb{R}\setminus \{x\}    
\end{equation*}
which is a contradiction, since $\mathbb{R}$ is not countably infinite and removing one point does not change this fact. Notice the intersection of the basis elements is solely $x$. That is because the topology is $T_1$, which means we can find some open set containing $x$ that separates $x$ and any other arbitrary point. Then by basis definition, there is some $B_n \subseteq O_x$.\\

This same $T_1$ condition idea with basis elements gives equivalence of a limit point and $\omega$-limit point in the next section. The key takeaway is that directions in proofs that aim to prove the existence of some topological object implies the existence of a sequential characterization require first countability.\\
\subsection{Nets}
We begin by introducing the notion of a directed set:
\begin{defn}
We define an ordering relation $\geq$ for the set $D$ and call it a directed set under the following conditions:
\begin{enumerate}
\item Reflexive: $\forall x \in D: x \geq x$
\item Transitive: $\forall x, y, z \in D: x \geq y \land y \geq z \implies x \geq z$
\item Directed: $\forall x, y \in D, \exists z \in D: (z \geq x \land z \geq y) \lor (x \geq z \land y \geq z)$
\end{enumerate}
\end{defn}
Notice that the first two properties align with a partial order. However a partial order also requires antisymmetry, which means:
\begin{equation*}
x \geq y \land y \geq x \implies x = y
\end{equation*}
which is not required for a directed set. This will become quite important later on. The third property (which is direction) tells us whether the net is upward directed (the left side condition) or downward directed (the right side), or both!\\

We also say that element $x$ follows element $y$ if $x \geq y$.\\

We view an equivalent property, for which WLOG, we assume an upward directed set $D$:
\begin{lem}
Given $(D, \geq)$ an upward directed set, then for all finite subsets $S$ of $D$:
\begin{equation*}
\forall S \subseteq D, \exists z \in D, \forall x \in S: z \geq x \iff \forall x, y \in D, \exists z \in D: z \geq x \land z \geq y 
\end{equation*}
In other words, all finite subsets of the directed set admit an upper bound.
\end{lem}
\begin{proof}
$\implies:$ The forward direction is trivial as a two element subset is finite.

$\impliedby:$ For the converse we prove via induction. The base case is a finite set of two elements, which we are given is true. Now assume it holds for a finite set of $n$ elements such that:
\begin{equation*}
S_n \coloneq \{x_i: 0 \leqslant i \leqslant n-1\} \land \forall x_i \in S_n, \exists z \in D: x_i \leq z
\end{equation*}
Now consider $S_{n+1} \coloneq S_n \cup \{x_n\}$ where $x_n$ is different from all elements in $S_n$. Then we get that, by the base case:
\begin{equation*}
\exists w \in D: x_n \leq w \land z \leq w
\end{equation*}
and finally by transitivity, we obtain that $\forall x_i \in S_{n+1}: x_i \leq w$. 
\end{proof} 

For the remainder of these notes, we assume a directed set $D$ is upward directed unless stated otherwise (the downward directed set notion will also be useful). Now a net is built on a directed set by the following:
\begin{defn}
A net is a function $\omega: D \xrightarrow{} X$ where $D$ is a directed set and $X$ is a topological space.  
\end{defn}
We define net convergence as the following. Now, to bring about the similarity to sequences, we shall denote a net as $\{x_\alpha\}_{\alpha \in D}$ where $D$ is the directed set. This is still a function where $x_\alpha = \omega(\alpha)$, but this representation shows $D$ as an indexing set.\\

Now in full generality of an arbitrary topological space, we say that the net converges or that $x_\alpha \xrightarrow{} x$ if:
\begin{equation*}
\forall U_x \in \uptau_X, \exists \alpha_0 \in D, \forall \alpha \geq \alpha_0: x_\alpha \in U_x
\end{equation*}
where $U_x$ is an arbitrary open neighborhood of the point $x \in X$ with the topology $\uptau_X$. For the majority of these notes, we shall work in metric spaces $(X, \varrho)$ (unless stated otherwise), so we can instead say:
\begin{equation*}
\forall \epsilon > 0, \exists \alpha_0 \in D, \forall \alpha \geq \alpha_0: x_\alpha \in B(c, \epsilon) \iff \varrho(x_\alpha, x) < \epsilon
\end{equation*}
so if our directed set $D = \mathbb{N}$, we obtain a sequence. We can alternatively characterize convergence by the following definition:
\begin{defn}
The tail of a net at $\alpha_0$ is defined as:
\begin{equation*}
T_{\alpha_0} \coloneq \{x_\alpha : \alpha \geq \alpha_0\}
\end{equation*}  
\end{defn}
Then a net converges to some $x$ if and only if $\forall \epsilon > 0, \exists \alpha_0 \in D$ such that:
\begin{equation*}
T_{\alpha_0} \subseteq B(x, \epsilon) 
\end{equation*}

\subsection{Cluster Points of Sequences}
We introduce a related concept from the lens of sequences:
\begin{defn}
A sequence frequently visits a set $A$ if:
\begin{equation*}
\forall n_0 \geqslant 0, \exists n \geqslant n_0: x_n \in A
\end{equation*}
\end{defn}
So a cluster point of a sequence is defined as the following:
\begin{defn}
$x$ is a cluster point of $\{x_n\}_{n \in \mathbb{N}}$ if $\{x_n\}_{n \in \mathbb{N}}$ frequently visits all open neighborhoods of $x$.
\end{defn}
The notion of a cluster point is something we will bring about when talking about nets and sequences. Considering metric spaces, we get the definition of a cluster point as:
\begin{equation*}
\forall \epsilon > 0, \forall n_0 \geqslant 0, \exists n \geqslant n_0: x_n \in B(x, \epsilon)
\end{equation*}
which permits $x_n = x$ unlike a \textbf{limit point of a set}, though there is an eventual connection. We view the following lemma:
\begin{lem}
$x$ is a cluster point of $\{x_n\}_{n \in \mathbb{N}} \in X^\mathbb{N}$ if and only if $\exists \{n_k\}_{k \in \mathbb{N}} \in \mathbb{N}^\mathbb{N}: x_{n_k} \xrightarrow{} x$  
\end{lem}
\begin{proof}
$\implies:$ If $x$ is a cluster point, then:
\begin{equation*}
\forall \epsilon > 0, \forall n_0 \geqslant 0, \exists n \geqslant n_0: x_n \in B(x, \epsilon)
\end{equation*}
We DO NOT require Axiom of Choice for this proof, as the definition of the cluster point provides us a method of construction. Define the following:
\begin{equation*}
\epsilon_{m} \coloneq \frac{1}{m+1}
\end{equation*}
Then, for $m = 0$, choose $n_0 = 0$, which gives some $n \geqslant 0$ such that $x_n \in B(x, \epsilon_0)$. Call this the first term of the subsequence as $n_0$ (different from the first $n_0$). \\

Then for $\epsilon_1$, we can find some $n \geqslant n_0 + 1$ such that $x_n \in B(x, \epsilon_1)$. Repeat this for all $m \in \mathbb{N}$, so that:
\begin{equation*}
\forall k \in \mathbb{N}: n_k < n_{k+1} \land \forall \epsilon > 0, \exists k_0 \geqslant 0, \forall k \geqslant k_0: |x_{n_k} - x| < \epsilon
\end{equation*}
because by the Archimedian property, we are guaranteed to find some $m \in \mathbb{N}$ such that:
\begin{equation*}
\frac{1}{m+1} \leqslant \epsilon \implies \exists n_m \geqslant 0, \forall n_k \geqslant n_m: |x_{n_k} - x| < \frac{1}{k+1} \leqslant \frac{1}{m+1} \leqslant \epsilon
\end{equation*}
which now shows convergence by Definition 0.2 and can also be phrased in terms of Definition 0.3.\\

$\impliedby:$ If there exists a convergent subsequence, then:
\begin{equation*}
\forall \epsilon > 0, \exists n_0 \geqslant 0, \forall n \geqslant n_0: x_n \in B(x, \epsilon)
\end{equation*}
This lemma shows us that the limit of a sequence is always a cluster point (not necessarily the converse) because convergence is a stricter condition. We see that:
\begin{equation*}
\forall \epsilon > 0, \forall n_0' \geqslant 0, \exists n \geqslant n_0': x_n \in B(x, \epsilon)
\end{equation*}
because if $n_0' < n_0$, then $n = n_0$ satisfies the condition of a cluset point. If $n_0' \geqslant n_0$, then simply choosing $n = n_0'$ satisifes the condition because by convergence $T_{n_0} \subseteq B(x, \epsilon)$. 
\end{proof}
Notice how (aside from the construction of the subsequence) did not use any properties specific to that of integers. Another thing is the limit of a convergent sequence is a cluster point, but the converse is not necessarily true. This gives rise to nets as we only require the properties of directed sets.

%We define an additional useful property:
% \begin{defn}
% $x$ is a limit-cluster point of $\{x_n\}_{n \in \mathbb{N}} \in X^\mathbb{N}$ if $\{x_n\}_{n \in \mathbb{N}}$ frequently visits all punctured open neighborhoods of $x$.
% \end{defn}
% In other words:
% \begin{equation*}
% \forall \epsilon > 0, \forall n_0 \geqslant 0, \exists n \geqslant n_0: x_n \in B(x, \epsilon) \setminus \{x\}
% \end{equation*}
% Then the following lemma (applying the same logic as Lemma 0.2) follows:
% \begin{lem}
% $x$ is a limit-cluster point of $\{x_n\}_{n \in \mathbb{N}} \in X^\mathbb{N}$ if and only if $\exists \{n_k\}_{k \in \mathbb{N}}: x_{n_k} \xrightarrow{} x \land (\forall k \in \mathbb{N}: x_{n_k} \neq x)$  
% \end{lem}
% Now in this case, even the limit of a convergent sequence may not be a limit-cluster point (consider a constant or eventually constant sequence). 
%some change%another one%and another one%and another one haha%
%allright i am making some changes let us see!%another one!%
%hey another edit%another one%and another one%and another one, ill see the github commit soon%another one!$
%alright final comment before the 5:30 commit!%

\subsection{Subsets of Directed Sets}
\begin{defn}
Given a directed set $D$, we call a subset $E \subseteq D$ cofinal if:
\begin{equation*}
\forall \alpha \in D, \exists \alpha' \in E: \alpha' \geq \alpha
\end{equation*}
\end{defn}
Cofinality is key in defining the concept of subnets (for which we will explore two separate definitions). We also see the following connection with cluster points/frequently in a set:
\begin{lem}
Given some topological space $X$ and some subset $A$, a net $\{x_\alpha\}_{\alpha \in D}$ in $X$ is frequently in $A$ if and only if there exists a cofinal subset $E$ such that:
\begin{equation*}
\forall \alpha \in E: x_\alpha \in A
\end{equation*}
\end{lem}
This follows from definition 0.6 given above. If there exists a cofinal subset $E$, then:
\begin{equation*}
\forall \alpha \in D, \exists \alpha' \in E: \alpha' \geq \alpha
\end{equation*}
Combined with the fact that the image of $E$ is in $A$, we see that:
\begin{equation*}
\forall \alpha \in D, \exists \alpha' \geq x: x_{\alpha'} \in A
\end{equation*}
This only proves one direction, but the other direction is constructive in nature (to show the existence of $E$ simply follows from the definition of the net being frequently in $A$).\\

We also see that cofinal subsets are directed. Let $E$ be a cofinal subset of a directed set $D$. Then by subset relations, it inherits the ordering ($\geq$) and is reflexive. We need to prove property 3 in Definition 0.1 only.\\

Consider some $\alpha', \alpha'' \in E$. We get that there must be some $\alpha \in D$ such that:
\begin{equation*}
\alpha \geq \alpha' \land \alpha \geq \alpha''
\end{equation*}
Now by cofinality, there exists some $\alpha''' \in E$ such that $\alpha''' \geq \alpha$, so by transitivity, the directnedness of $E$ follows.\\

We say that a net converges if it is eventually in all open neighborhoods of $x$.\\

Let us examine the primary definition (\textbf{General Topology, Kelley}) we will use for a subnet:
\begin{defn}
Given a net $\{x_\alpha\}_{\alpha \in D}$ a subnet is the function $\omega \circ f$, where $f: E \xrightarrow{} D$ and $\omega$ for some directed set $(E, \succcurlyeq)$ was the initial net such that $\omega: D \xrightarrow{} X$.\\

We denote the subnet by $\{x_{f(\beta)}\}_{\beta \in E}$, which must adhere by the following property:
\begin{equation*}
\forall \alpha \in D, \exists \beta_0 \in E, \forall \beta \succcurlyeq \beta_0: f(\beta) \geq \alpha
\end{equation*}
\end{defn}
The intuition behind this definition\\

Now consider some cofinal subset of $D$, which we will call $E$ with the same ordering $\geq$. Then cofinality gives us that:
\begin{equation*}
\forall \alpha \in D, \exists \beta_0 \in E: \beta_0 \geq \alpha \iff \forall \alpha \in D, \exists \beta_0 \in E, \forall \beta \geq \beta_0: \beta \geq \alpha
\end{equation*}
which follows from transitivty. So $\{x_\alpha\}_{\alpha \in E}$ is a subnet of the initial net with $f$ being the identity function, which is similar to convergence properties of subsequences and cluster points. Using this fact, we can extend Lemma 0.3 to the following:
\begin{lem}
Given some topological space $X$ and some subset $A$, a net $\{x_\alpha\}_{\alpha \in D}$ in $X$ is frequently in $A$ if and only if there exists a subnet contained in $A$.   
\end{lem}
\begin{proof}
The proof utilizes the results of the previous definition. Lemma 0.3 gave us equivalence with the existence of some cofinal subset $E \subseteq D$ that maps into $A$. A cofinal subset is essentially a subnet under composition with the identity function, thus giving our desired equivalence.\\
\end{proof}

This can be extended to the notion of convergence, essentially bringing about a parallel lemma to Lemma 0.2 stated in terms of nets:
\begin{lem}
$x \in X$ is a cluster point of $\{x_\alpha\}_{\alpha \in D}$ if and only if $\exists \{x_{f(\alpha')}\}_{\alpha' \in E}: x_{f(\alpha')} \xrightarrow{} x$
\end{lem}
%this is quite an interesting proof that I need to examine properly and rigorously%
%though I think I have enough information to expound upon nets in analysis and rigorously define such instances%
\begin{proof}
A cluster point can be defined as one for which the net is frequently in all neighborhoods of the point $U_x$. By the previous lemma, we can say that $x$ is a cluster point if and only if there is a subnet contained in each neighborhood, but that is quite weak and difficult to check. Let us define $E$ as the following:
\begin{equation*}
E \coloneq \{(\alpha', U_x) : x_{\alpha'} \in U_x\}
\end{equation*}
$\implies$: We show that $E$ is a directed set that brings about a subnet $\{x_{f(\alpha')}\}_{\alpha' \in E}$. $E$ is directed under the product ordering whereby:
\begin{equation*}
(\alpha', U_x) \succcurlyeq (\alpha_1', V_x) \iff \alpha_1' \geq \alpha' \land V_x \subseteq U_x
\end{equation*}
So taking any two $(\alpha', U_x)$ and $(\alpha_1', V_x)$ in $E$, we see that by the ordering of $D$:
\begin{equation*}
\exists \alpha_2' \in D: \alpha_2' \geq \alpha' \land \alpha_2' \geq \alpha_1'
\end{equation*}
and similarly for the neighborhoods, we have that:
\begin{equation*}
V_x \cap U_x \subseteq V_x \land V_x \cap U_x \subseteq U_x
\end{equation*}
Now the pair $(\alpha_2', V_x \cap U_x)$ may not belong to $E$, but by the cluster point definition for all neighborhoods of $x$, we have the existence of some $\alpha_3'$, such that:
\begin{equation*}
x_{\alpha_3'} \in V_x \cap U_x \implies (\alpha_3', V_x \cap U_x) \in E
\end{equation*}
Now that we have established that $E$ is directed (the other two properties trivially follow from subset inclusion and directedness of $D$). Now, we define the function $f: E \xrightarrow{} D$ such that:
\begin{equation*}
\forall (\alpha', U_x) \in E: f(\alpha', U_x) = \alpha'
\end{equation*}
We can show that this is a subnet (specifically a Kelley subnet, but later we will show it is also a Willard subnet). The existence of the cluster point $x$ gives us that $D$ is frequently in all neighborhoods:
\begin{equation*}
\forall U_x, \forall \alpha \in D, \exists \alpha_1 \geq \alpha: x_{\alpha_1} \in U_x \implies (\alpha_1, U_x) \in E
\end{equation*} 
\begin{equation*}
\implies \forall U_x, \forall \alpha \in D, \exists (\alpha_1, U_x) \in E: f(\alpha_1, U_x) \geq \alpha
\end{equation*}
And by directedness of $E$, we see that:
\begin{equation*}
\implies \forall U_x, \forall \alpha \in D, \exists (\alpha_1, U_x) \in E, \forall (\alpha_2, V_x) \succcurlyeq (\alpha_1, U_x): f(\alpha_2, V_x) \geq \alpha
\end{equation*}
Thus proving that $\{x_{f(\alpha, V_x)}\}_{(\alpha, V_x) \in E}$ is a Kelley subnet. It follows that the subnet converges to $x$ as:
\begin{equation*}
\implies \forall U_x, \exists (\alpha_1, U_x) \in E, \forall (\alpha_2, V_x) \succcurlyeq (\alpha_1, U_x): x_{\alpha_2} \in V_x \subseteq U_x
\end{equation*}
as it is eventually in each neighborhood of $x$. We now prove the converse:

$\impliedby:$ If $x$ is not a cluster point of the net, then:
\begin{equation*}
\exists U_x, \exists \alpha \in D, \forall \alpha_1 \geq \alpha: x_{\alpha_1} \notin U_x
\end{equation*}
Then for all possible subnets/cofinal subsets:
\begin{equation*}
\forall \alpha \in D, \exists \alpha_0' \in E, \forall \alpha' \geq \alpha_0': f(\alpha') \geq \alpha
\end{equation*}
we get that:
\begin{equation*}
\implies \exists U_x, \exists \alpha_0' \in E, \forall \alpha' \geq \alpha_0': f(\alpha') \in D \land f(\alpha') \geq \alpha \implies x_{f(\alpha')} \notin U_x
\end{equation*}
which means that any subnet is not eventually in a neighborhood of $x$, thus there exists no subnet that converges to $x$.
\end{proof}
Notice we can do the same proof in a first countable space with sequences; however, this construction will not give us a subsequence (because it is not a map from $\mathbb{N}$ to $\mathbb{N}$), rather a subnet, showing that a subsequence is a subnet but not necessarily the converse.\\

Let us examine a different definition of a subnet (\textbf{General Topology, Willard}):
\begin{defn}
Given a net $\{x_\alpha\}_{\alpha \in D}$, we define a subnet $\{x_{f(\alpha')}\}_{\alpha' \in E}$ by the following conditions:
\begin{enumerate}
\item $\forall \alpha \in D, \exists \alpha_0' \in E, f(\alpha_0') \geq \alpha$
\item $f: E \xrightarrow{} D$ is order preserving such that:
\begin{equation*}
\forall \alpha_0', \alpha_1' \in E: \alpha_0' \succcurlyeq \alpha_1' \implies f(\alpha_0') \geq f(\alpha_1')
\end{equation*}
\end{enumerate}   
\end{defn}
The Willard definition implies the Kelley definition (we can use condition two to derive the Kelley condition), however the converse is not necessarily true.\\

The Willard definition is strictly stronger with the order preserving function, but the intuition behind both is that the indexing terms $f(E)$ go "further" in $D$. Our traditional definition of a subsequence is a Willard subnet, but it can possibly be relaxed to fit Kelley's definition.\\

We write such a condition as follows:
\begin{defn}
We define a (Willard) subsequence as $\{x_{n_k}\}_{k \in \mathbb{N}}$ where $\{n_k\}_{k \in \mathbb{N}}$ is a strictly increasing sequence. In other words, we have a function:
\begin{equation*}
\exists f: \mathbb{N} \xrightarrow{} \mathbb{N} \ \text{such that } \forall k, j \in \mathbb{N}: k > j \implies f(k) = n_k > n_j = f(j)
\end{equation*}
The first condition is satisfied the condition that $\forall k \in \mathbb{N}: n_k \geqslant k$, which can be proven inductively.\\

We define a (Kelley) subsequence as $\{x_{n_k}\}_{k \in \mathbb{N}}$ where for $\{n_k\}_{k \in \mathbb{N}}$, we have a that:
\begin{equation*}
\forall n \in \mathbb{N}, \exists k_0 \in \mathbb{N}, \forall k > k_0: n_k > n
\end{equation*}
Note that Kelley does use $\geqslant$ in this definition, however this effectively means a singular value of a single index can create a subsequence. 
\end{defn}
What this means is we can have a Kelley subsequence with indices such as $n_0 = 1, n_1 = 3, n_2 = 2$, while this would not be allowed in a Willard subsequence.\\

Now in terms of analysis, because of the total ordering of $\mathbb{R}$, we will deal with Willard subnets naturally, though I plan to explicitly use this terminology defined in this section regarding nets in the next section. That way, the connection between the concepts and nets is stronger.

\subsection{Nets In Analysis}
\subsubsection{Oscillation as a Net}
We have the oscillation characterization of continuity of a function $f$ at a point $x_0 \in \text{int}(\text{Dom}(f))$ as:
\begin{equation*}
\forall \epsilon > 0, \exists \delta_0 > 0, \forall \delta \leqslant \delta_0: \sup_{x, y \in B(x_0, \delta)}|f(x) - f(y)| \leqslant \epsilon
\end{equation*}
where we write the net as $\{\text{osc}(f, B(x_0, \delta))\}_{\delta > 0}$ and:
\begin{equation*}
\text{osc}(f, B(x_0, \delta)) \coloneq \sup_{x, y \in B(x_0, \delta)}|f(x) - f(y)|
\end{equation*}
where the downward directed set is $(\mathbb{R}_{> 0}, \leqslant)$ as the domain of the net.\\

Let us view an important related result:
\begin{thm}
$f: \mathbb{R} \xrightarrow{} \mathbb{R}$ is continuous at $x_0$ if and only if for all $\delta$-nets:
\begin{equation*}
\forall \{x_\delta\}_{\delta > 0}: x_\delta \xrightarrow{} x_0 \implies f(x_\delta) \xrightarrow{} f(x_0)
\end{equation*}
\end{thm}
\begin{proof}
$\implies:$ If $f$ is continuous at $x_0$, then we have that:
\begin{equation*}
\forall \epsilon > 0, \exists \delta_0 > 0, \forall x \in B(x_0, \delta_0): f(x) \in B(f(x_0), \epsilon)
\end{equation*}
Then consider all $\delta$-nets that converge to $x_0$, we have that:
\begin{equation*}
\forall U_{x_0}, \exists \delta_\alpha > 0, \forall \delta \leqslant \delta_\alpha: x_\delta \in B(x_0, \delta_\alpha) \subseteq U_{x_0}
\end{equation*}
Then letting $U_x = B(x_0, \delta_0)$, for all convergent $\delta$-nets:
\begin{equation*}
\exists \delta_\alpha > 0, \forall \delta \leqslant \delta_\alpha: x_\delta \in B(x_0, \delta_0)
\end{equation*}
which by continuity implies that $f(x_\delta) \in B(f(x_0), \epsilon)$, so for all convergent $\delta$-nets:
\begin{equation*}
\implies \forall \epsilon > 0, \exists \delta_\alpha > 0, \forall \delta \leqslant \delta_\alpha: f(x_\delta) \in B(f(x_0), \epsilon)
\end{equation*}
proving that $\forall \{x_\delta\}_{\delta > 0}: f(x_\delta) \xrightarrow{} f(x_0)$.\\

$\impliedby:$ We prove by contrapositive. If $f$ is not continuous at $x_0$, then we have that:
\begin{equation*}
\exists \epsilon > 0, \forall \delta > 0, \exists x \in B(x_0, \delta): f(x_\delta) \notin B(f(x_0), \epsilon)
\end{equation*}
and we want to show there exists a convergent $\delta$-net whose image is not converging to $f(x_0)$. For this we require Axiom of Choice (and notice we do not require first countability).\\

Construct the net $\{x_\delta\}_{\delta > 0}$ such that:
\begin{equation*}
\forall \delta > 0: x_\delta \in B(x_0, \delta) \land f(x) \notin B(f(x_0), \epsilon)
\end{equation*}
which is clearly a convergent net. However the image net $\{f(x_\delta)\}_{\delta > 0}$ is not convergent because:
\begin{equation*}
\exists \epsilon > 0, \forall \delta > 0, \exists \delta_0 = \delta: f(x_{\delta_0}) \notin B(f(x_0), \epsilon)
\end{equation*}
thus proving the converse. 
\end{proof}
What this theorem says is that a function $f$ is (pointwise) continuous at $x_0$ if and only if for all nets convergent to $x_0$, their image converges to $f(x_0)$. This idea allows us to formalize the following notion:
\begin{equation*}
\lim_{x_\delta \xrightarrow{} x_0} f(x_\delta) = f(x_0)
\end{equation*}
where below the limit is the convergence of the $\delta$-net and the entire operation represents convergence of the image. A similar theorem can be stated for differentiability and as we will see in the future, integrability.\\

While in general, texts define the above notation as the limit via $\epsilon-\delta$ and punctured $\delta$-balls, to maintain consistency, we will write a net converges as $x_\delta \xrightarrow{} x_0$ and use the above notation only when a convergence of a net implies the convergence of its image under a function.\\

Now relative to this theorem, we can call oscillation as the "worst case $\delta$-net", which if it converges to zero, then all other nets must too converge to the function value (follows from completeness). We can manipulate the oscillation criterion further:
\begin{equation*}
\forall \epsilon > 0, \exists \delta_0 > 0, \forall \delta \leqslant \delta_0: - \epsilon \leqslant 0 \leqslant \sup_{x, y \in B(x_0, \delta)}|f(x) - f(y)| \leqslant \epsilon  
\end{equation*}
\begin{equation*}
\implies \forall \epsilon > 0, \exists \delta_0 > 0, \forall \delta \leqslant \delta_0: \left|\text{osc}(f, B(x, \delta))-0\right| \leqslant \epsilon
\end{equation*}
so $T_{\delta_0} \subseteq B(0, \epsilon)$, thus proving the net:
\begin{equation*}
\text{osc}(f, B(x_0, \delta)) \xrightarrow{} 0
\end{equation*}
The main goal of this section is to jump to nets for integration and sequences of functions. However, as I write notes down for those, I will definitely revisit and improve this section. The reason why I wanted to highlight oscillation as a net is so the reader can draw similar parallels with integration and function criterions.
\pagebreak

\section{Compactness}
In the previous section, we discussed notions of convergence and adherence for sequences. We now discuss these notions for sets. Let us view the following theorem:

\subsection{Limit Points}
We start off with a topological proof, which will be key to discussing limit and adherent points:
\begin{lem}
The closure of a set $A$ is equal to all of its adherent points. 
\end{lem}
\begin{proof}
$\implies:$ We prove by contrapositive. If $x$ is not an adherent point, then:
\begin{equation*}
\exists \epsilon > 0: B(x, \epsilon) \cap A = \varnothing \implies B(x, \epsilon) \subseteq X \setminus A
\end{equation*} 
So $x$ is an interior point of $X \setminus A$, which means:
\begin{equation*}
x \in \text{int}(X \setminus A) = X \setminus \overline{A}
\end{equation*} 

$\impliedby:$ We once again prove by contrapositive. If $x \notin \overline{A}$, then:
\begin{equation*}
x \in X \setminus \overline{A} \subseteq X \setminus A \implies x \in X \setminus A
\end{equation*}
Since $x \setminus \overline{A}$ is open, then $x$ must be an interior point, so:
\begin{equation*}
\exists \epsilon > 0: B(x, r) \subseteq X \setminus A \implies \exists \epsilon > 0: B(x, r) \cap A = \varnothing
\end{equation*}
\end{proof}
The limits of sequences ARE NOT necessarily limit points. Given some subset $E \subseteq \mathbb{R}$, $x$ is a limit point of $E$ if:
\begin{equation*}
\forall \epsilon > 0: B(x, \epsilon) \cap E \neq \varnothing
\end{equation*}
Notice how this only requires the intersection to be nonempty; which means there could possibly be a finite number of points in the intersection (but we will see this is not possible). Let us look at the following lemma:
\begin{lem}
$x$ is a limit point of $E$ if and only if
\begin{equation*}
\exists \{x_n\}_{n \in \mathbb{N}} \in (E \setminus \{x\})^\mathbb{N}: x_n \xrightarrow{} x
\end{equation*}
\end{lem}
\begin{proof}
$\implies$: If $x$ is a limit point, then using the first countability of metric spaces, consider:
\begin{equation*}
\epsilon_n \coloneq \frac{1}{n+1}
\end{equation*} 
Then by axiom of choice, choose $x_n \in B(x, \epsilon_n) \cap E \setminus \{x\}$, so that:
\begin{equation*}
\forall n \in \mathbb{N}: |x_n - x| < \frac{1}{n+1} \implies x_n \xrightarrow{} x
\end{equation*}

$\impliedby$: Now for the converse. If there exists such a sequence $\{x_n\}_{n \in \mathbb{N}}$ such that:
\begin{equation*}
\forall \epsilon > 0, \exists n_0 \geqslant 0, \forall n \geqslant n_0: |x_n - x| < \epsilon
\end{equation*}
then we see that
\begin{equation*}
\forall \epsilon > 0: B(x, \epsilon) \cap E \setminus \{x\} \neq \varnothing
\end{equation*}
as $\forall n \geqslant n_0: x_n \in B(x, \epsilon)$ and by definition of the sequence $x_n \in E \setminus \{x\}$. Thus proving equivalence.
\end{proof}

Now we introduce the following notion:
\begin{defn}
An $\omega$-limit point of a set $E$ is one for which there exists a countably infinite number of points of $E$ within each punctured $\epsilon$-ball around $x$.
\end{defn}
It follows that every limit point in $\mathbb{R}$ is an $\omega$-limit point from Lemma 1.2, the sequential definition. This is because:
\begin{equation*}
\forall \epsilon > 0, \exists n_0 \geqslant 0, \forall n \geqslant n_0: x_n \in B(x, \epsilon) \cap E \setminus \{x\}
\end{equation*}
and the natural numbers $n_0$ and beyond are countably infinite (an injection of $\mathbb{N}$ into this set can easily be shown), thus $x$ is an $\omega$-accumulation point. The weakest sufficient condition is a $T_1$ space for such equivalence.\\

Now if the set $E$ consists of some sequence $\{x_n\}_{n \in \mathbb{N}}$, then $E \coloneq \{x_n : n \in \mathbb{N}\}$, we see that $x$ is a limit point of $E$ if and only if there exists a convergent \textbf{subsequence} converging to $x$. This is not immediately trivial as one of the conditions we require for a subsequence is that:
\begin{equation*}
\forall \{n_k\}_{k \in \mathbb{N}}: n_k < n_{k+1s}
\end{equation*}
%test edit%
Observe the definition of sequential closure of some subset $E$:
\begin{defn}
$\overline{E}$ is the closure of $E$ if and only if
\begin{equation*}
\forall \{x_n\}_{n \in \mathbb{N}} \in E^\mathbb{N}: x_n \xrightarrow{} x \implies x \in \overline{E}
\end{equation*}
\end{defn}
So by the equivalence stated in Lemma 1.2, $\{x_n\}_{n \in \mathbb{N}}$ belongs to $E^\mathbb{N}$ as well (it is not guaranteed that $x \in E$), but by sequential closure $x$ or the limit point must belong in $\overline{E}$.\\

Now steering our attention back to sequences, the limit of a sequence may not be a limit point. Consider an eventually constant or constant sequence. We see that it converges to whatever value it is eventually constant at; but there is an empty $\epsilon$-punctured ball intersection. So a limit point is a limit of a specific type of sequence, but not vice versa.\\

This distinction will be quite crucial in the following section.

\subsection{Limit Points and Sequences}
In the previous section, we introduced the notion of a cluster point of a sequence, for which we proved:
\begin{equation*}
x \ \text{is a cluster point of } \{x_n\}_{n \in \mathbb{N}} \in X^\mathbb{N} \iff \exists \{n_k\}_{k \in \mathbb{N}} \in \mathbb{N}^\mathbb{N}: x_{n_k} \xrightarrow{} x
\end{equation*}
but nothing prevented $x_{n_k}$ from being perhaps a constant sequence such that all terms equal $x$ itself. Now this section focused on limit points of \textbf{sets}. So consider the set:
\begin{equation*}
S \coloneq \{x_n : n \in \mathbb{N}\}
\end{equation*}
which is the set of all possible values that appear in the sequence. Notice that the sequence allows for the repetition of values, while $S$ does not. So if $S$ is finite such that:
\begin{equation*}
S = \{x_0, ..., x_m\}
\end{equation*}
we notice that because a metric space is $T_1$, each of these points can be isolated from eachother. Thus, $S$ has no limit point. Now assume $S$ were countabily infinite, which means the sequence has infinitely many distinct terms. Then:
\begin{lem}
\begin{equation*}
x \ \text{is a limit point of } S \iff \exists \{n_k\}_{k \in \mathbb{N}} \in \mathbb{N}^\mathbb{N}: x_{n_k} \xrightarrow{} x \land (\forall k \in \mathbb{N}: x_{n_k} \neq x)
\end{equation*}
\end{lem}
This aligns with Lemma 1.2 with the \textbf{extra condition} that it must be a subsequence, which implies:
\begin{equation*}
n_k \xrightarrow{} \infty \land \forall k \in \mathbb{N}: n_k < n_{k+1}
\end{equation*}
\begin{proof}
$\implies:$ Define $\epsilon_m \coloneq 1/(m+1)$. Then:
\begin{equation*}
\forall k \in \mathbb{N}: B(x, \epsilon_k) \cap S \setminus \{x\} \neq \varnothing
\end{equation*}
Then we have some $x_n \in S$ that belongs to the ball $B(x, \epsilon_k)$. Define this as the $n_k$ term of the "sub"-sequence. Now, this does not necessarily guarantee that the sequence we constructed $\{n_k\}_{k \in \mathbb{N}}$ is necessarily strictly increasing.\\

By Sequential Bolzano Weierstrass, we notice that the indices are bounded below by $0$ and that any sequence of these indicies must contain either a strictly increasing subsequence or a non-increasing subsequence. Let us define the notion of dips in the sequence as:
\begin{equation*}
D \coloneq \{m \in \mathbb{N}: (\forall j > m: n_j > n_m)\}
\end{equation*}
In the case that $D$ is finite, we cannot find a strictly increasing subsequence, but we could possibly find a non-increasing subsequence. So, consider:
\begin{equation*}
m_0 \coloneq \sup (D) + 1
\end{equation*}
so that the peaks are no longer included in the subsequence.\\ 

Then for all $m \geqslant m_0$, $m \notin D$, which tells us that:
\begin{equation*}
\exists j > m: n_j \leqslant n_m
\end{equation*}
so we choose $m_{k+1} \coloneq \sup\{j > m_k : n_{m_j} \leqslant n_{m_k}\}$, which is possible because we have an upper bound. So, the initial value $n_{m_0}$ acts as an upper bound and $0$ is a lower bound for the indices.\\

This means either the sequence is finite (which is a contradiction), or that the sequence is constant/eventually constant. What this means is that the resultant subsequence $\{x_{n_{m_k}}\}_{k \in \mathbb{N}}$ takes a finite number of values.\\

Because the space is $T_1$ (this condition is necessary for the proof), then we can find a neighborhood that contains $x$ and none of the values in the sequence, which too is a contradiction. Thus we must be able to find a strictly increasing subsequence of indices, such that:
\begin{equation*}
\forall k \in \mathbb{N}: n_{m_k} < n_{m_{k+1}}
\end{equation*} 
and since $x_{n_k} \xrightarrow{} x$, we have that $x_{n_{m_k}} \xrightarrow{} x$ and satisfies the conditions for a subsequence of $\{x_n\}_{n \in \mathbb{N}}$. 
%have to work on this more%

$\impliedby:$ If there exists a convergent subsequence, then:
\begin{equation*}
\forall \epsilon > 0, \exists k_0 \geqslant 0, \forall k \geqslant k_0: x_{n_k} \in B(x, \epsilon)
\end{equation*}
Now given all terms of the subsequence do not equal $x$ and $S$ is countably infinite, we see that:
\begin{equation*}
\forall \epsilon > 0: B(x, \epsilon) \cap S \setminus \{x\} \neq \varnothing
\end{equation*}
and furthermore $x$ is an $\omega$-limit point. 
\end{proof}
This lemma is the engine behind the equivalence of limit point compactness and sequential compactness (the forward direction utilizes first countability).\\

Going forward, we shall separate the terminologies so that we talk about cluster points of sequences/nets and limit/adherent points of sets. We provide a lemma for their equivalence below:
% \begin{lem}
% If $S \coloneq \{x_n : n \in \mathbb{N}\}$ is countably infinite for some $\{x_n\}_{n \in \mathbb{N}} \in X^\mathbb{N}$, then:
% \begin{equation*}
% x \ \text{is a limit point of } S \iff x \ \text{is a limit-cluster point of } \{x_n\}_{n \in \mathbb{N}}
% \end{equation*}
% \end{lem}
% \begin{proof}
% The proof follows from the equivalent condition proposed in Lemma 0.3. $x$ is a limit-cluster point of a sequence if and only if there exists a subsequence that converges to $x$ such that none of its terms equal $x$.\\

% This is the same equivalent condition proposed in Lemma 1.3, thus proving the above lemma.
% \end{proof}

Another lemma requiring a stronger sufficient condition is the following:
\begin{lem}
If $\{x_n\}_{n \in \mathbb{N}} \in X^\mathbb{N}$ consists of all DISTINCT terms (each term appears at most once), then:
\begin{equation*}
x \ \text{is a limit point of } S \iff x \ \text{is a cluster point of } \{x_n\}_{n \in \mathbb{N}}  
\end{equation*} 
\end{lem}
\begin{proof}
Lemma 0.2 provides us with the equivalence that $x$ is a cluster point of the sequence if and only if there exists a convergent subsequence to $x$. Now since the sequence itself has all distinct points, so does the convergent subsequence.\\

One of the terms could perhaps be $x$, but since it appears at most once, omission of a finite number of points does not change convergence of the sequence. We prove this simply below. Let $x$ be the $n_j$-th element of the original subsequence derived from the definition of cluster point. Then, for the new subsequence with $x$ omitted:
\begin{equation*}
\forall \epsilon > 0, \exists k_0' \coloneq \max\{k_0, j+1\}, \forall k \geqslant k_0': |x_{n_k} - x| < \epsilon
\end{equation*}
where $k_0$ comes from the convergence of the subsequence containing $x$. Thus we can always find a tail of the new sequence within the $\epsilon$-neighborhood of $x$. If $k_0 < j$, then $k_0' = j+1$, and by convergence of the initial subsequence, the new one certainly converges as:
\begin{equation*}
T'_{k_0'} = T_{j+1} \subseteq T_{k_0} \subseteq B(x, \epsilon)
\end{equation*}
In the case that $k_0 > j$, we see that:
\begin{equation*}
T'_{k_0'} = T_{k_0} \subseteq B(x, \epsilon)
\end{equation*}
\end{proof}
Lemma 1.4 and 1.5 will be key when discussing sequential compactness. Notice that for this proof, we used the sequential characterizations to draw equivalence; however, it can be done without first countability. The converse can be done in general topological spaces and follows from the definition of a cluster point.\\

The forward direction, requires $T_1$. We prove by contradiction. Assume $x$ is a limit point (thus an $\omega$-limit point), but not a cluster point, such that:
\begin{equation*}
\exists U_x \subseteq X, \exists n_0 \geqslant 0, \forall n \geqslant n_0: x_n \notin U_x
\end{equation*}
Then there are $n_0$ values indexed by the values $0,...n_0 - 1$ that may belong to $U_x$. One of those values could be $x$ (but at most one, because $S$ contains all distinct values), which contradicts the fact that $x$ is an $\omega$-limit point of $S$. While Theorem 1.3 deals with equivalence in $\mathbb{R}$ or general metric spaces (which are both $T_1$ and first countable), we assume the equivalence.\\

The reader will notice that Theorem 1.4 expands the previous theorem to more generality, and this distinction we have mentioned will be key to explaining the implications and how we can still have equivalence in non-first countable yet $T_1$ spaces.

\subsection{Sequential Compactness}
\begin{thm}[Sequential Bolzano Weierstrass]
Any bounded sequence has a convergent subsequence.
\end{thm}
To prove this fact, we show first that any sequence in $\mathbb{R}$ has a monotonic subsequence. 
\begin{proof}
Let us define the following set, we call "peaks" of a sequence:
\begin{equation*}
P \coloneq \{n \in \mathbb{N}: (\forall j > n: x_j < x_n)\}
\end{equation*}
In the case that $P$ is countably infinite, then we set:
\begin{equation*}
n_0 \coloneq \inf P \land \forall k > 0: n_k \coloneq \inf \left(P \setminus \{n_0,...,n_{k-1}\}\right)
\end{equation*}
In which case we are guaranteed that $\forall k \in \mathbb{N}: n_{k+1} > n_k$ which tells us that $x_{n_{k+1}} < x_{n_k}$, showing that there exists a strictly decreasing subsequence. In the case that $P$ is finite, we choose the following:
\begin{equation*}
n_0 \coloneq \sup(P) + 1
\end{equation*}
that way the peaks are fully ommited from the subsequence, as they will prove to be problematic in the construction. Then for any $k \in \mathbb{N}$, we are guaranteed that $n_k \notin P$ (since $\{n_k\}_{k \in \mathbb{N}}$ is strictly increasing and $n_0 \notin P$), which means:
\begin{equation*}
\exists j > n_k: x_j \geqslant x_{n_k}
\end{equation*}
So we simple define:
\begin{equation*}
\forall k \in \mathbb{N}: n_{k+1} \coloneq \inf\{j > n_k: x_j \geqslant x_{n_k}\}
\end{equation*}
in which case we obtain a non-decreasing subsequence. Now any bounded sequence has a monotone subsequence, which converges. What this means is given some sequence $\{x_n\}_{n \in \mathbb{N}}$, we have that:
\begin{equation*}
\exists y \in \mathbb{R}, \exists r > 0, \forall n \in \mathbb{N}: x_n \in B(y, r)
\end{equation*}
so the sequence has a convergent subsequence within the ball, but such a subsequence may not converge within the ball. This is due to the topological idea that the limit point certainly belongs in the closure of the set but not necessarily in the open set. Consider the bounded sequence:
\begin{equation*}
\{x_n\}_{n \in \mathbb{N}} \in (0, 1)^\mathbb{N}: \left(\forall n \in \mathbb{N}: x_n \coloneq \frac{1}{n+2}\right)
\end{equation*}
Then $x_n \xrightarrow{} 0$ yet $0 \notin (0, 1)$. This idea of closedness plays a key role in sequential compactness. 
\end{proof}

\begin{defn}
A sequentially compact interval of $\mathbb{R}$ is one for which every sequence has a convergent subsequence that converges to a point within the interval. 
\end{defn}
The previous Theorem 1.1 leads to a key theorem:
\begin{thm}
An interval of $\mathbb{R}$ is sequentially compact if and only if it is closed and bounded. 
\end{thm}
$\impliedby$: This direction follows from Sequential Bolzano Weierstrass. If the interval is bounded, then any sequence has a convergent monotone subsequence. Since it is closed, the limit of the sequence, which is an adherent point, belongs in the closure.\\

$\implies$: Consider any convergent sequence $\{x_n\}_{n \in \mathbb{N}}$ such that $\forall n \in \mathbb{N}: x_n \in I$. Then by sequential compactness:
\begin{equation*}
\exists \{n_k\}_{k \in \mathbb{N}}: x_{n_k} \xrightarrow{} x \land x \in I
\end{equation*}
which forces $x_n \xrightarrow{} x$. Thus all convergent sequences converge to points within the set (which by sequential closure) means $I$ is closed. Now assume the interval is not bounded, in which case:
\begin{equation*}
\forall r > 0, \forall y \in \mathbb{R}: I \nsubseteq B(y, r)
\end{equation*}
Or equivalently (to bring about the sequential format, we use first countability):
\begin{equation*}
\forall n \in \mathbb{N}, \forall y \in \mathbb{R}, \exists x \in I: x \notin B(y, n+1)
\end{equation*}
Choose some arbitrary $y \in \mathbb{R}$. Then $\forall n \in \mathbb{N}$ choose $x_n$ such that $x_n \notin B(y, n+1)$, so that:
\begin{equation*}
\forall n \in \mathbb{N}: |x_n - y| > n+1 > 1/2
\end{equation*}
Now assume there existed some $z$ such that some subsequence indexed by $\{n_k\}_{k \in \mathbb{N}}$ converged to $z$. Then:
\begin{equation*}
\forall k \in \mathbb{N}: |x_{n_k} - y| \leqslant |x_{n_k} - z| + |y-z|
\end{equation*}
\begin{equation*}
\implies \forall k \in \mathbb{N}: ||x_{n_k} - y| - |y-z|| \leqslant |x_{n_k} - z|
\end{equation*}
Taking limit $k \xrightarrow{} \infty$ we see that $n_k \xrightarrow{} \infty$, thus:
\begin{equation*}
||x_{n_k} - y| - |y-z|| \xrightarrow{} 0
\end{equation*}
which is a contradiction because $|x_{n_k} - y| > n_k + 1$ which tells us that $|x_{n_k} - y| \xrightarrow{} \infty$, so $x_{n_k} \nrightarrow z$ and since $z$ is arbitrary, this holds true for all such $z \in \mathbb{R}$. 

\begin{lem}
For any set $I \subseteq \mathbb{R}$, $I$ is bounded if and only if it is totally bounded (ONLY for $\mathbb{R}$).
\end{lem}
\begin{proof}
$\impliedby$: If $I$ is bounded, then:
\begin{equation*}
\exists r > 0, \exists y \in \mathbb{R}: I \subseteq B(y, r) \iff \exists r > 0, \exists y \in \mathbb{R}, \forall x \in I: x \in B(y, r)
\end{equation*}
We aim to show $I$ is totally bounded, which means:
\begin{equation*}
\forall \epsilon > 0, \exists x_0, ..., x_m \in I: I \subseteq \bigcup_{i=0}^{m}B(x_i, \epsilon)
\end{equation*}
This converse direction is necessarily true; however, the forward direction is not necessarily true in all spaces. Consider the following:
\begin{equation*}
M \coloneq \max\{|x_i - x_j| : 0 \leqslant i, j \leqslant m\}
\end{equation*}
Then, choose some $x_i$; so without loss of generality, we choose $x_0$:
\begin{equation*}
\forall x \in I, \exists x_i: x \in B(x_i, \epsilon) \implies |x-x_0| \leqslant |x-x_i| + |x_i - x_0| < \epsilon + M
\end{equation*}

$\implies:$ Now to prove the forward direction. We have that $I$ is bounded, which means:
\begin{equation*}
\exists y \in \mathbb{R}, \exists r > 0: I \subseteq B(y, r)
\end{equation*}
Now we can show that $B(y, r)$ is totally bounded. Simply consider the uniform partition:
\begin{equation*}
\Pi_B \coloneq \left\{(y-r) + \frac{2r}{n}i: 0 \leqslant i \leqslant n\right\}
\end{equation*}
whereby:
\begin{equation*}
\exists n \in \mathbb{N}: \frac{2r}{n} < \frac{\epsilon}{2}
\end{equation*}
\begin{equation*}
\implies \forall \epsilon > 0, \exists \Pi_B: I \subseteq B(y, r) \subseteq \bigcup_{x_i \in \Pi_B}B(x_i, \epsilon/2)
\end{equation*}
Now $\forall x_i \in \Pi_B$ where $B(x_i, \epsilon/2) \cap I \neq \varnothing$, choose $x_i'$ such that:
\begin{equation*}
x_i' \in B(x_i, \epsilon/2) \cap I
\end{equation*}
and add it to the partition $\Pi_I$. Then we observe that:
\begin{equation*}
\forall x \in I, \exists x_i \in \Pi_B: x \in B(x_i, \epsilon/2)
\end{equation*}
\begin{equation*}
\implies \exists x_i' \in \Pi_I: |x - x_i'| \leqslant |x - x_i| + |x_i - x_i'| < \epsilon/2 + \epsilon/2 = \epsilon
\end{equation*}
thus, $I$ is totally bounded such that,
\begin{equation*}
\forall \epsilon > 0, \exists \Pi_I: I \subseteq \bigcup_{x_i' \in \Pi_I}B(x_i', \epsilon)
\end{equation*}
\end{proof}
Now the main assumption we made here is that every open ball in $\mathbb{R}$ is totally bounded (which we can show via uniform partitions). The rest of the argument is not unique to $\mathbb{R}$. This can give insight into what spaces obey this Heine-Borel compactness property.

\subsubsection{Bounded/Totally-Bounded Equivalence}
NOTE: I will look into this later, it has something to do with finite dimensions/functional analysis. The main thing to investigate is in what spaces are open balls totally bounded.

\subsection{Limit Point Compactness}
We now discuss an equivalent form of compactness called limit point compactness (which due to the first countability of $\mathbb{R}$) is essentially equivalent to the sequential characterization.

\begin{defn}
A set $K \subseteq \mathbb{R}$ is limit point compact if for all infinite (countably or uncountably) sets $E \subseteq K$, there exists a limit point of $E$ in $K$.
\end{defn}

\begin{thm}
A set $K$ is limit point compact if and only if it is sequentially compact in $\mathbb{R}$.
\end{thm}
$\implies:$ If $K$ is limit point compact, then every infinite subset is limit point compact. Then consider each sequence $\{x_n\}_{n \in \mathbb{N}} \in K^\mathbb{N}$ and define:
\begin{equation*}
S \coloneq \{x_n : n \in \mathbb{N}\}
\end{equation*}
If $S$ is finite, then by pigeonhole principle, we can find a constant convergent subsequence. If $S$ is countably infinite, then there exists some point $x \in K$ which is a limit point of $S$. By Lemma 1.3, we have the existence of a subsequence with all terms being different from $x$ converging to $x$.\\

$\impliedby:$ Consider every infinite subset of $K$. From that we can always find a countably infinite subset. Let us define this subset as $E$. Then consider an arbitrary sequence $\{x_n\}_{n \in \mathbb{N}} \in E^\mathbb{N}$, such that $S \coloneq \{x_n : n \in \mathbb{N}\}$ is countably infinite and all terms of the sequence are unique (this is possible because $E$ is countably infinite).\\

Then, the sequence must admit a convergent subsequence, which means it admits a cluster point. By Lemma 1.4, we see that this cluster point must be a limit point of $S \subseteq E$, thereby a limit point of $E$.
% \begin{proof}
% $\implies$: Now consider any sequence $\{x_n\}_{n \in \mathbb{N}} \in K^\mathbb{N}$, which is an infinite subset, so by limit point compactness it admits a limit point. Now:
% \begin{equation*}
% \forall n \in \mathbb{N}: \{x_n\}_{n \in \mathbb{N}} \cap B\left(x, \frac{1}{n+1}\right) \neq \varnothing
% \end{equation*}
% So to construct the converging subsequence (note this requires the Axiom of Choice), we choose:
% \begin{equation*}
% \forall k \in \mathbb{N}: x_{n_k} \in \{x_n\}_{n \in \mathbb{N}} \cap B\left(x, \frac{1}{k+1}\right)
% \end{equation*}
% such that $\forall k \in \mathbb{N}: n_{k+1} > n_k$. (\textbf{Not so sure about this proof because can we justify such a choice whereby $n_{k+1} > n_k$}).\\

% An alternative prove is by contrapositive. Consider $K$ is not sequentially compact, then:
% \begin{equation*}
% \exists \{x_n\}_{n \in \mathbb{N}} \in K^\mathbb{N}, \forall \{n_k\}_{k \in \mathbb{N}}, \forall k \in \mathbb{N}: n_k < n_{k+1} \implies \forall x \in K: x_{n_k} \nrightarrow x
% \end{equation*}
% which means there exists a sequence with no convergent subsequence. In that case, the sequence is a subsequence of itself, which means there exists an infinite subset of $K$ with no limit point, which follows from:
% \begin{equation*}
% \exists x \in K: x \ \text{is a limit point of } \{x_n\}_{n \in \mathbb{N}} \iff \exists \{n_k\}_{k \in \mathbb{N}}: n_k < n_{k+1} \land x_{n_k} \xrightarrow{} x
% \end{equation*}
% the proof of which is provided in the section about nets.\\

% $\impliedby:$ Consider any infinite subset of $K$. Then taking a countable infinite subset of that, we obtain a sequence $\{x_n\}_{n \in \mathbb{N}} \in K^\mathbb{N}$ of DISTINCT points.\\

% By sequential compactness, there exists a convergent subsequence converging to a point $x$. Notice that we do not want one value repeating infinitely many times; otherwise we may not satisfy sequential characterization of limit points.\\

% Now it is possible that $x$ appears as a value, but at most once within the subsequence. Ommission of $x$ from this subsequence still results in a convergent subsequence. Assume $x$ appears at some index $n_j$, then:
% \begin{equation*}
% \forall \epsilon > 0, \exists k_0 \geqslant 0, \forall k \geqslant k_0: |x_{n_k} - x| < \epsilon
% \end{equation*}
% so if $j \geqslant k_0$, we still have infinite values other than $x$ from the subsequence within the $\epsilon$-ball around $x$. This is trivial if $j < k_0$, thus $x$ is an $\omega$-limit point or equivalently a limit point.
% \end{proof}
% Notice that the $\implies$ direction is not always true in general topological spaces as we used the first countability of metric spaces. The equivalence really stems from the idea of nets/countable nets or sequences. Limit point compactness is identical to sequential compactness in first countable spaces; the only difference being the former describes net-wise behavior, while the latter describes sequential.\\

% But the equivalence fundamentally comes from the equivalent characterization of limit points. 

We now see that first countability can be avoided by weakening the notion of sequential compactness to cluster-point compactness:
\begin{thm}
A set $K$ is limit point compact if and only if it is cluster-point compact in a $T_1$ space.
\end{thm}
\begin{proof}
$\implies$: This direction requires $T_1$. For all sequences, if $S \coloneq \{x_n : n \in \mathbb{N}\}$ is finite, then there is a convergent subsequence by pigeonhole principle, which is a cluster point.\\

If $S$ is countably infinite (then it admits a limit point) and $\{x_n\}_{n \in \mathbb{N}}$ has no constant subsequence (elements appear finitely many times), then take the subsequence omitting repetitions (resulting in a subsequence with entirely unique elements).\\

We utilize the alternative proof of Lemma 1.4, showing that in a $T_1$ space, if $x$ is a limit point of $S$, then it is a cluster point of the sequence.\\

$\impliedby$: This direction requires neither $T_1$ nor first countability. Consider any countably infinite subset $E \subseteq K$ and enumerate it by the naturals, giving a sequence of distinct terms. Then if $x$ is a cluster point of the sequence, we have that:
\begin{equation*}
\forall U_x, \forall n_0 \geqslant 0, \exists n \geqslant n_0: x_n \in U_x
\end{equation*}
\begin{equation*}
\implies \forall U_x, \exists n \geqslant 0: x_n \in U_x \implies \forall U_x: U_x \cap E \setminus \{x\} \neq \varnothing
\end{equation*}
which follows from the definitions as well as the fact that the sequence consists of distinct terms. Thus, $x$ is a limit point of $E$ (not necessarily an $\omega$-limit point, but this was not required by our given definition).
\end{proof}
\pagebreak

\section{Topological Compactness}
Open covers give us a way to talk about topological compactness in a highly geometric fashion. The two previous notions, limit point and sequential compactness are quite difficult to intuitively understand geometrically.
\begin{defn}
Given some subset $E$ of a topological space $X$, an open cover of $E$ is a collection:
\begin{equation*}
\{O_\alpha: \alpha \in I \land O_\alpha \in \uptau_X\}
\end{equation*}
such that:
\begin{equation*}
E \subseteq \bigcup_{\alpha \in I} O_\alpha
\end{equation*}
\end{defn}
We say that a subset $E$ of a topological space $X$ (which can include $X$ itself) is topologically compact if EVERY open cover has a finite subcover. By the Heine-Borel definition in $\mathbb{R}$, we see that $(0, 1)$ is not compact. For this, define the open cover:
\begin{equation*}
\left\{\left(0, 1-\frac{1}{n+1}\right) : n \in \mathbb{N}\right\}
\end{equation*}
which admits no finite subcover. Consider any finite subcollection of this open cover. We see that its union must be some interval $(0, 1-1/(m+1))$ for some $m \in \mathbb{N}$ (the largest $m$). However, we are guaranteed to always find some real number $r$ such that:
\begin{equation*}
1 - \frac{1}{m+1} < r < 1
\end{equation*}
which tells us that there cannot be any finite subcover.\\

Notice that if a metric space is topologically compact, then it must be totally bounded, as consider the following cover for some fixed $\epsilon > 0$ and $E \subseteq X$:
\begin{equation*}
\{B(x, \epsilon): x \in E\}
\end{equation*}
Since $X$ is topologically compact, it admits a finite subcover, which means:
\begin{equation*}
\exists x_0,...,x_m \in E: E \subseteq \bigcup_{0 \leqslant i \leqslant m}B(x_i, \epsilon)
\end{equation*}
The converse is not necessarily true, as evident by our example with $(0, 1)$. Because it is bounded, by Lemma 1.5, it is totally bounded (this is true only for spaces where open balls are totally bounded (possibly?)), yet we showed it is neither Heine-Borel compact, nor topologically compact.

\subsection{Topological and Limit Point Compactness}
Interestingly enough, these two notions are not equivalent in the general sense. We only have the following implication:
\begin{thm}
Given a general topological space $X$, if $X$ is topologically compact, then it is limit point compact.
\end{thm}
\begin{proof}
We prove by contrapositive. First, if $X$ is a finite space, then every open cover must be finite, and it is vacuously limit point compact (because all infinite subsets have a limit point, of which there are no infinite subsets). So we deal with infinite topological spaces (those with infinite underlying points).\\

If $X$ is not limit point compact, then taking an infinite subset $E$:
\begin{equation*}
\exists E \subseteq X, \forall x \in X, \exists U_x: E \cap U_x \setminus \{x\} = \varnothing
\end{equation*}
Now take the collection of all such $U_x$ as $\mathcal{U}_x \coloneq \{U_x : E \cap U_x \setminus \{x\} = \varnothing\}$ as an open cover (let us assume that each $U_x$ is an open neighborhood). Then taking any finite subcover at points $\{x_0, ..., x_m\}$, we see that each $x_i$ is either in $E$ or not in $E$.\\

In either case, we see that not all of $E$ is covered by the finite subcover, since removing finite number of $x_i$ points from $E$ does not change the fact that it is infinite, which means there will always be points of $E$ not covered by any finite subcover, and since $E \subseteq X$, this holds for $X$ as well.\\

So, we have found an open cover $\mathcal{U}_x$ with no finite subcover, proving $X$ is not topologically compact.
\end{proof}


\pagebreak

\section{Integration and Nets}
The definition of Riemann and Riemann-Stieltjes integrability are fundamentally based on nets of partitions. We examine both varieties of integrals here as well as total variation.

\subsection{Mesh of Partitions}
Given a compact interval $[a, b] \subseteq \mathbb{R}$, we define a partition as:
%I need to better define what partitions are to define orderings on them%
%I think I should take a look at the Henstock-Kurtzweil Integral%
%Perhaps the concept of delta-fine partitions will help with this%
\begin{equation*}
\Pi_{[a, b]} \coloneq \{t_i : 0 \leqslant i \leqslant n\}
\end{equation*}
such that $t_0 = a$ and $t_n = b$ and:
\begin{equation*}
\forall 0 \leqslant i \leqslant n-1: t_i < t_{i+1}
\end{equation*}
so that:
\begin{equation*}
a < t_1 < \dots < t_{n-1} < b
\end{equation*}
We then define the mesh of a partition $\Pi$ as the following:
\begin{equation*}
||\Pi|| \coloneq \max_{0 \leqslant i \leqslant n-1}|t_{i+1} - t_i|
\end{equation*}
We furthermore say $\Pi_1$ is a refinement of or is finer than $\Pi_0$ if $\Pi_0 \subseteq \Pi_1$. Due to the uniqueness of the supremum, it is guaranteed that the mesh relation is a function. The following lemma will prove to be quite useful:
\begin{lem}
If $\Pi_1 \supseteq \Pi_0$, then $||\Pi_1|| \leqslant ||\Pi_0||$.
\end{lem}
\begin{proof}
Consider the points $\{x_0, ..., x_m\}$ to be added to the partition $\Pi_0$ resulting in $\Pi_1$, such that:
\begin{equation*}
\Pi_0 \coloneq \{t_i: 0 \leqslant i \leqslant n\}
\end{equation*}
\begin{equation*}
\Pi_1 \coloneq \{u_j : 0 \leqslant j \leqslant n+m\}
\end{equation*}
such that $u_0 = a$ and $u_{n+m} = b$. We further define:
\begin{equation*}
\forall 0 \leqslant i \leqslant n-1: T_i \coloneq \{[u_j, u_{j+1}] : 0 \leqslant j \leqslant n+m - 1 \land [u_j, u_{j+1}] \subseteq [t_i, t_{i+1}]\}   
\end{equation*}
Now we use basic properties of outer measure where $A \subseteq B \implies |A| \leqslant |B|$, which leads to the following:
\begin{equation*}
\forall 0 \leqslant i \leqslant n-1, \forall [u_j, u_{j+1}] \in T_i: |[u_j, u_{j+1}]| \leqslant |[t_i, t_{i+1}]|
\end{equation*}
Since this is true for all $i \in [0, n-1]$, we see that:
\begin{equation*}
\forall 0 \leqslant i \leqslant n-1, \forall [u_j, u_{j+1}] \in T_i: |[u_j, u_{j+1}]| \leqslant ||\Pi_0||
\end{equation*}
\begin{equation*}
\implies \forall 0 \leqslant j \leqslant n+m-1: |[u_j, u_{j+1}]| \leqslant |[t_i, t_{i+1}]|| \leqslant ||\Pi_0||
\end{equation*}
and because we have an upper bound on the measure of each interval, we get that $||\Pi_1|| \leqslant ||\Pi_0||$. 
\end{proof}
When we refer to a partition in the next subsection, we will worked with marked partitions containing sampling points in each interval $t_i^*$. So when we say $\Pi$ is a partition of $[a, b]$, we mean that it is the ordered pair:
\begin{equation*}
\Pi \coloneq (\{t_i\}_{i \in [0, n]}, \{t_i^*\}_{i \in [1, n]})
\end{equation*}
Then, we say that $\Pi_1 \supseteq \Pi_0$ if and only if $\{t_i\}_{i \in [0, n]}^{\Pi_1} \supseteq \{t_i\}_{i \in [0, m]}^{\Pi_0}$. Note that this is not antisymmetric in the sense that a bi-directional relation does not imply the tagged partitions are the same. It only implies the underlying endpoints are the same. (This ordering relation independent of marked points is as per Apostol's formulation.)\\


We make this distinction as it ensures that the net we introduce is a function (otherwise one partition may have two different Riemann-Stieltjes sums) due to different marked points. The following lemmas can be proved even when considering the sequence of tags for the ordering relation; however, it is easier when solely considering the underlying endpoints.

\subsection{Riemann-Stieltjes Integrability}
Given two functions $f, g \in \mathbb{R}^\mathbb{R}$, we call $f$ the integrand and $g$ the integrator.
\begin{defn}
The function $f$ is Riemann-Stieltjes integrable on $[a, b]$ (with respect to $g$) if:
\begin{equation*}
\forall \epsilon > 0, \exists \delta_0 > 0, \forall \Pi: ||\Pi|| < \delta_0 \implies \left|S(f, g, \Pi) - \int_{a}^{b}fdg\right| < \epsilon
\end{equation*}
\end{defn}
Let us now examine how nets play a role here. Let $P$ be the set of all finite partitions of $[a, b] \subseteq \mathbb{R}$. We say that $\Pi_1 \leq \Pi_0$ iff $||\Pi_1|| \leqslant ||\Pi_0||$.\\

Thus $(P, \leq)$ is a directed set, because:
\begin{equation*}
\forall \Pi_0, \Pi_1 \in P: \Pi_1, \Pi_0 \subseteq \Pi_1 \cup \Pi_0 \implies \Pi_1 \cup \Pi_0 \leq \Pi_1, \Pi_0
\end{equation*}
by Lemma 3.1, showing $P$ is downward directed (the other properties follow from the mesh function mapping to $\mathbb{R}$). The net we will define will then be the net of Stieltjes sums $\{S(f, g, \Pi)\}_{\Pi \in P}$:
\begin{equation*}
\forall \Pi \in P: S(f, g, \Pi) \coloneq \sum_{i=1}^{n}f(t_i^*)(t_i - t_{i-1})
\end{equation*}
For brevity, we may also use the following notation for Stieltjes sums:
\begin{equation*}
S(f, g, \Pi) = \sum_{T_i \in \Pi}f(t_i^*)\Delta T_i
\end{equation*}
where each $T_i \coloneq [t_{i-1}, t_i]$. Such a net converges if:
\begin{equation*}
\forall \epsilon > 0, \exists \Pi_\epsilon \in P, \forall \Pi \leq \Pi_\epsilon: S(f, g, \Pi) \in B\left(\int_{a}^{b}fdg, \epsilon \right)
\end{equation*}

\begin{lem}
The conditions in Definition 3.1 and the convergence of $\{S(f, g, \Pi)\}_{\Pi \in P}$ are equivalent.
\end{lem}
\begin{proof}
$\implies:$ If there exists some $\delta_0 > 0$, then we can take the uniform net $\Pi_\epsilon$ with mesh at most $\frac{\delta_0}{2}$. By the Archimedian property:
\begin{equation*}
\frac{\delta_0}{2(b-a)} > 0 \implies \exists n \in \mathbb{N}: 1 \leqslant n\frac{\delta_0}{2(b-a)} \implies \frac{b-a}{n} \leqslant \frac{\delta_0}{2} 
\end{equation*}
So we can define $\Pi_\epsilon$ as the following:
\begin{equation*}
\Pi_\epsilon \coloneq \left\{a + i \cdot \frac{b-a}{n} : 0 \leqslant i \leqslant n\right\}
\end{equation*}
Then:
\begin{equation*}
\forall \epsilon > 0, \exists \Pi_\epsilon \in P, \forall \Pi \in P: \Pi \leq \Pi_\epsilon \implies ||\Pi|| \leqslant ||\Pi_\epsilon|| \leqslant \frac{\delta_0}{2} \leqslant \delta_0
\end{equation*}
and by the criterion in Definition 3.1:
\begin{equation*}
\implies S(f, g, \Pi) \in B\left(\int_{a}^{b}fdg, \epsilon\right)
\end{equation*}
showing that the net $S(f, g, \Pi) \xrightarrow{} \int_{a}^{b}fdg$.\\

$\impliedby:$ We now prove the converse. We have the convergence of our Stieltjes sum net:
\begin{equation*}
\forall \epsilon > 0, \exists \Pi_\epsilon \in P, \forall \Pi \leq \Pi_\epsilon: S(f, g, \Pi) \in B\left(\int_{a}^{b}fdg, \epsilon\right)
\end{equation*}
Now take $\delta_0 \coloneq ||\Pi_\epsilon||/2$, so we get:
\begin{equation*}
\forall \epsilon > 0, \exists \delta_0 > 0, \forall \Pi \in P: ||\Pi|| < \delta_0 \implies ||\Pi|| < ||\Pi_\epsilon||/2 < ||\Pi_\epsilon||
\end{equation*}
\begin{equation*}
\implies \Pi \leq \Pi_\epsilon \implies S(f, g, \Pi) \in B\left(\int_{a}^{b}fdg, \epsilon\right)
\end{equation*}
\end{proof}

\subsubsection{Alternative Integrability Criteria}
Notice the generality in the net $\{S(f, g, \Pi)\}_{\Pi \in P}$ that we constructed. Lemma 3.1 tells us that if $\Pi_1 \supseteq \Pi_0$, then in this case $\Pi_1 \leq \Pi_0$ or equivalently $||\Pi_1|| \leqslant ||\Pi_0||$.\\

The converse however is not necessarily true. If $\Pi_1 \leq \Pi_0$ that does not necessarily mean that $\Pi_1 \supseteq \Pi_0$. We now look at a more specific definition:
\begin{defn}
The function $f$ is Riemann-Stieltjes integrable on $[a, b]$ with respect to $g$ if
\begin{equation*}
\forall \epsilon > 0, \exists \Pi_\epsilon \in P, \forall \Pi \supseteq \Pi_\epsilon: S(f, g, \Pi) \in B\left(\int_{a}^{b}fdg, \epsilon\right)
\end{equation*}
\end{defn}
We can "net-ify" this definition by observing the following construction:
\begin{equation*}
P_\epsilon \coloneq \{\Pi \in P : \Pi \supseteq \Pi_\epsilon\} 
\end{equation*}
We can show that $\forall \epsilon > 0$, the $P_\epsilon$ is a cofinal subset of $P$ such that:
\begin{equation*}
\forall \Pi_0 \in P, \exists \Pi \in P_\epsilon: \Pi \leq \Pi_0
\end{equation*}
For any arbitrary $\Pi_0 \in P$, we see that $\Pi_0 \cup \Pi_\epsilon \supseteq \Pi_\epsilon$ and by Lemma 3.1:
\begin{equation*}
\Pi_0 \cup \Pi_\epsilon \supseteq \Pi_0 \implies \Pi_0 \cup \Pi_\epsilon \leq \Pi_0 \land \Pi_0\cup \Pi_\epsilon \in P_\epsilon
\end{equation*}
Since the image of $P_\epsilon$ under the identity is cofinal, it is the domain for a Kelley/Willard Subnet (either classification holds because the identity is order preserving).\\

Thus Definition 3.2 can be restated as $f$ is RS-integrable with respect to $g$ if:
\begin{equation*}
\forall \epsilon > 0, \exists \{S(f, g, \Pi)\}_{\Pi \in P_\epsilon}: S(f, g, \Pi) \in B\left(\int_{a}^{b}fdg, \epsilon\right)
\end{equation*}
which to put into words means, for all neighborhoods of $\int_{a}^{b}fdg$, there is a subnet contained in that neighborhood. By Lemma 0.5, this means $\int_{a}^{b}fdg$ is a cluster point of the net $\{S(f, g, \Pi)\}_{\Pi \in P}$.
%this seems to be related to cluster points%
%I think it implies the value of the integral is a cluster point because it implies that for all neighborhoods, there is a subnet contained within it%
%which means the net is frequently in all neighborhoods, thus a cluster point%
%so perhaps to disprove RS-integrability, a criterion could be if:%
%\exists \epsilon > 0, \exists \Pi_\epsilon \in P, \forall \Pi \leq \Pi_\epsilon: S(f, g, \Pi)\notin B(L, \epsilon)%
%yes this would make sense because then this implies the quantifier switch of 2nd and 3rd quantifiers%
%which leads us to the negation of the condition equivalent in Lemma 3.2%


\begin{thm}
If $f$ is Riemann-Stieltjes integrable with respect to $g$ by the criterion in Definition 3.1, it is Riemann-Stieltjes integrable by the criterion in Definition 3.2.
\end{thm}
\begin{proof}
The criterion in Definition 3.1 states in terms of nets that for the net $\{S(f, g, \Pi)\}_{\Pi \in P}$:
\begin{equation*}
S(f, g, \Pi) \xrightarrow{} \int_{a}^{b}fdg
\end{equation*}
so the entire net converges to the value of the integral, which means any subnet must converge to the value of the integral as well, thereby implying Definition 3.2's criterion. To put it more succinctly, the limit of a net is a cluster point of a net, but not vice versa.
\end{proof}
Another characterization equivalent to Definition 3.2 would be the following:
\begin{equation*}
\forall \epsilon > 0, \forall \Pi_0 \in P, \exists \Pi \leq \Pi_0: S(f, g, \Pi) \in B\left(\int_{a}^{b}fdg, \epsilon\right)
\end{equation*}
A major topic we shall analyze has now emerged. In the Riemann integral (when $g$ is the identity function), we actually have equivalence of these notions; however, in the Stieltjes integral, we do not always. So the goal now is to figure out some sufficient conditions for equivalence.

Going forward, we will use the following two notions:
\begin{equation*}
(1): \ \forall \epsilon > 0, \exists \Pi_\epsilon \in P, \forall \Pi \leq \Pi_\epsilon: S(f, g, \Pi) \in B\left(\int_{a}^{b}fdg, \epsilon\right)
\end{equation*}
\begin{equation*}
(2): \ \forall \epsilon > 0, \exists \Pi_\epsilon \in P, \forall \Pi \supseteq \Pi_\epsilon: S(f, g, \Pi) \in B\left(\int_{a}^{b}fdg, \epsilon\right)
\end{equation*}

We also require the following lemma:
\begin{lem}
If $f$ is Riemann integrable by criterion (2), then it is bounded on $[a, b]$. 
\end{lem}
\begin{proof}
So Riemann Integrability gives us that:
\begin{equation*}
\forall \epsilon > 0, \exists \Pi_\epsilon \in P, \forall \Pi \supseteq \Pi_\epsilon: R(f, \Pi) \in B\left(\int_{a}^{b}fdx, \epsilon\right)
\end{equation*}
where $R(f, \Pi) = S(f, \text{id}, \Pi)$ such that $\text{id}(x) = x$ (the identity function). Then, for some arbitary $\epsilon = 1$, we see that:
\begin{equation*}
\forall \Pi \supseteq \Pi_\epsilon: \left|R(f, \Pi) - \int_{a}^{b}fdx\right| < 1
\end{equation*}
\begin{equation*}
\implies \forall \Pi \supseteq \Pi_\epsilon: |R(f, \Pi)| < \left|\int_{a}^{b}fdx\right| + 1
\end{equation*}
Now for any arbitrary $x \in [a, b]$, if $x$ is already a marked point of $\Pi_\epsilon$ in some interval $[t_{j-1}, t_j]$, then:
\begin{equation*}
f(x)(t_j - t_{j-1}) = R(f, \Pi_\epsilon) - \sum_{i \in [1, n]\setminus\{j\}}f(t_i^*)(t_i - t_{i-1})
\end{equation*}
\begin{equation*}
\implies |f(x)|(t_j - t_{j-1})| \leqslant |R(f, \Pi_\epsilon)| + \sum_{i \in [1, n]\setminus\{j\}}|f(t_i^*)| \cdot ||\Pi_\epsilon||
\end{equation*}
Let $T$ be the length of the smallest interval in $\Pi_\epsilon$, then:
\begin{equation*}
\implies |f(x)| \leqslant \frac{1}{T}\left(\left|\int_{a}^{b}fdx\right| + 1|\right) + \frac{||\Pi_\epsilon||}{T}\sum_{i\in [1, n]\setminus\{j\}}|f(t_i^*)|
\end{equation*}
Now if $x$ is not a marked point of $\Pi_\epsilon$ (it may be in an interval or the endpoint of an interval), then assume $x \in [t_{j-1}, t_j]$ in $\Pi_\epsilon$. Then define $\Pi$ such that $t_j^* = x$ in $\Pi$, with the remaining endpoints and marked points the same as $\Pi_\epsilon$.\\

Since $\Pi$ retains the same interval endpoints, we can say that $\Pi \supseteq \Pi_\epsilon$ and $||\Pi|| \leqslant ||\Pi_\epsilon||$ and that $T_\Pi = T$. Thus:
\begin{equation*}
|f(x)|(t_j - t_{j-1})| \leqslant |R(f, \Pi)| + \sum_{i \in [1, n]\setminus\{j\}}|f(t_i^*)| \cdot ||\Pi|| 
\end{equation*}
\begin{equation*}
|f(x)| \leqslant \frac{1}{T}\left(\left|\int_{a}^{b}fdx\right| + 1\right) + \frac{||\Pi_\epsilon||}{T}\sum_{i \in [1, n]\setminus\{j\}}|f(t_i^*)|
\end{equation*}
Since $x$ was arbitrary, this bound holds for all $x \in [a, b]$ no matter where $x$ falls with respect to the marked points of $\Pi_\epsilon$. Thus $f$ is bounded, and by least upper bound property, we can write the supremum of $f$ as $||f||$. 
\end{proof}
Because criterion 1 implies criterion 2, we also have the following implication:
\begin{lem}
If $f$ is Riemann integrable by criterion (1), then it is bounded on $[a, b]$.
\end{lem}

\medskip
We can now prove the equivalence of the two criterions for Riemann integrable functions:
\begin{thm}
If $f$ is Riemann integrable, then:
\begin{equation*}
\forall \epsilon > 0, \exists \Pi_\epsilon \in P, \forall \Pi \leq \Pi_\epsilon: R(f, \Pi) \in B\left(\int_{a}^{b}fdx, \epsilon\right)
\end{equation*}
\begin{equation*}
\iff \forall \epsilon > 0, \exists \Pi_\epsilon \in P, \forall \Pi \supseteq \Pi_\epsilon: R(f, \Pi) \in B\left(\int_{a}^{b}fdx, \epsilon\right)
\end{equation*}
\end{thm}
\begin{proof}
$\implies$: This direction holds in full generality for Riemann-Stieltjes integrals, which we proved in the previous subsection.\\

$\impliedby:$ Given the second criterion:
\begin{equation*}
\forall \epsilon > 0, \exists \Pi_\epsilon \in P, \forall \Pi \supseteq \Pi_\epsilon: R(f, \Pi) \in B\left(\int_{a}^{b}fdx, \frac{\epsilon}{2}\right)
\end{equation*}
we aim to prove the first. Let $\Pi_\epsilon^{(1)}$ be the partition for the first criterion, defined as a uniform partition with mesh less than $\frac{\epsilon}{2(||f||\cdot m + 1)}$. Then we want to show:
\begin{equation*}
\forall \Pi \leq \Pi_\epsilon^{(1)}: R(f, \Pi) \in B\left(\int_{a}^{b}fdx, \epsilon\right)
\end{equation*}
By triangle inequality, we see the following:
\begin{equation*}
\forall \Pi \leq \Pi_\epsilon, \exists \Pi \cup \Pi_\epsilon \in P: \left|R(f, \Pi) - \int_{a}^{b}fdx\right| \leqslant \left|R(f, \Pi) - R(f, \Pi \cup \Pi_\epsilon)\right| + \left|R(f, \Pi \cup \Pi_\epsilon) - \int_{a}^{b}fdx\right|
\end{equation*}
\begin{equation*}
\leqslant |R(f, \Pi) - R(f, \Pi \cup \Pi_\epsilon)| + \frac{\epsilon}{2}
\end{equation*}
which follows from the fact that $\Pi \cup \Pi_\epsilon \supseteq \Pi_\epsilon$. Since $\Pi \cup \Pi_\epsilon$ is merely an intermediary partition, we can assign convenient marked points to bring about the $\epsilon$-estimate.\\



Consider the marked partition $\Pi$ with arbitrary marked points. Now to its unmarked version, add one point in $\Pi_\epsilon$ let us call it $e \in [t_{j-1}, t_j]$. Then let $e$ be the marked point of $[t_{j-1}, e]$ and $[e, t_{j}]$. For the other intervals, utilize the marked points from $\Pi$ and call this partition $\Pi_1$. Then:
\begin{equation*}
R(f, \Pi) - R(f, \Pi_1) = f(t_j^*)(t_j - t_{j-1}) - f(e)(t_j - t_{j-1}) = (f(t_j^*) - f(e))(t_j - t_{j-1})
\end{equation*}
Take another point from $\Pi_\epsilon$ and add it to $\Pi_1$ in the same fashion. If no other point of $\Pi_\epsilon$ is in that interval, then the above estimate holds. Otherwise, consider some $h \in [e, t_j]$. Then:
\begin{equation*}
R(f, \Pi_1) - R(f, \Pi_2) = (f(e) - f(h))(t_j - e)
\end{equation*}
where the marked point of the subintervals $[e, h], [h, t_j]$ is now $h$. Repeat this process for all $m$ points in $\Pi_\epsilon$. In the remaining subintervals without a marked point, choose that of $\Pi$ since it means these intervals do not contain any points of $\Pi_\epsilon$. Now observe:
\begin{equation*}
R(f, \Pi) - R(f, \Pi_2) = (f(t_j^*) - f(e))(t_j - t{j-1}) + (f(e) - f(h))(t_j - e) \leqslant 2||f|| \cdot 2 \cdot ||\Pi||
\end{equation*}
Since the process was done $m$ times and our mesh of $\Pi$ is bounded, we get the following estimate:
\begin{equation*}
|R(f, \Pi) - R(f, \Pi \cup \Pi_\epsilon)| + \frac{\epsilon}{2}\leqslant 2||f|| \cdot m \cdot \frac{\epsilon}{2(||f|| \cdot m + 1)} + \frac{\epsilon}{2} \leqslant \epsilon
\end{equation*}
What allows us to carry out this proof is the fact that the second Riemann sum is based on a partition $\Pi \cup \Pi_\epsilon \supseteq \Pi$ and allows it to work for any arbitrary choice of marked points for $\Pi$ as per our definition. 
\end{proof}

To uncover the assumption that makes this proof work, we do the same proof except for the criterion equivalent to criterion 1:
\begin{equation*}
\forall \epsilon > 0, \exists \delta_0 > 0, \forall \Pi \in P: ||\Pi|| < \delta_0 \implies R(f, \Pi) \in B\left(\int_{a}^{b}fdx, \epsilon\right)  
\end{equation*}
Now the question is, how do we obtain a $\delta_0$? The trick we used to bring about our $\epsilon$ estimate in the previous proof was the control the mesh of the partition $\Pi_\epsilon^{(1)}$. So if we say:
\begin{equation*}
0 < \delta_0 \leqslant \frac{\epsilon}{2(||f|| \cdot m + 1)}
\end{equation*}
then we get that $\forall \Pi \in P : ||\Pi|| < \delta_0$, then:
\begin{equation*}
\Pi \cup \Pi_\epsilon \supset \Pi \implies ||\Pi \cup \Pi_\epsilon|| \leqslant ||\Pi|| < \delta_0 \leqslant \frac{\epsilon}{2(||f||\cdot m + 1)}
\end{equation*}
Then the proof is the same as Theorem 3.2 once this relation is established. The underlying assumption we use here however is uniform continuity of the identity function $(\text{id})$. We can prove that $(\text{id})$ is continuous (simply by taking $\epsilon = \delta$).\\

Thus on a compact interval $[a, b]$, it must be uniformly continuous such that:
\begin{equation*}
\forall \epsilon > 0, \exists \delta_0 > 0, \forall x, y \in [a, b]: |x-y| < \delta_0 \implies |\text{id}(x) - \text{id}(y)| = |x-y| < \epsilon
\end{equation*}
We actually see that it is $1$-Lipschitz because $|\text{id}(x) - \text{id}(y)| \leqslant |x-y|$, so we can set $\delta_0 = \epsilon$. Thus:
\begin{equation*}
\delta_0 \leqslant \frac{\epsilon}{2(||f||\cdot m + 1)} \implies |\text{id}(x) - \text{id}(y)| < \frac{\epsilon}{2(||f|| \cdot m + 1)}
\end{equation*}
but now what if the function were replaced with an arbitrary function $g(x)$. Then, we would require uniform continuity of $g$, so this is a peek into the sufficient conditions for equivalence of the two criteria in Riemann-Stieltjes integration.\\

Notice also that in the process of finding the difference between $R(f, \Pi)$ and $R(f, \Pi \cup \Pi_\epsilon)$, we obtained terms such as $(f(t_j^*) - f(e))(t_j - t_{j-1})$, we have a difference of terms of $f$. If we can control the distance between $t_j^*, e$ for example and are given that $f$ is continuous, then we could bring about an $\epsilon$-estimate.\\

In order to examine these two ideas, we need to introduce the topics of bounded variation and Darboux Formulations/Nets.

\subsection{Total Variation}
We define the total variation of a function $f$ over the compact interval $[a, b]$ as the following:
\begin{equation*}
V(f, [a, b]) \coloneq \sup_{\Pi \in P}\sum_{i=1}^{n}|f(t_i) - f(t_{i-1})|
\end{equation*}
where each $\Pi$ is an unmarked partition. We say that $f$ is of bounded variation on $[a, b]$ if $\exists C \in \mathbb{R}: V(f, [a, b]) = C < \infty$ or in shorthand $V(f, [a, b]) < \infty$.\\

Monotonic (without loss of generality, we assume non-decreasing) functions are of bounded variation, since:
\begin{equation*}
\forall \Pi \in P: \sum_{i=1}^{n}|f(t_i) - f(t_{i-1})| = f(b) - f(a)
\end{equation*}
and so by supremum properties, we have that $V(f, [a, b]) = f(b) - f(a)$. Another useful lemma is the following:

% \begin{lem}
% A function $f$ is continuously differentiable if and only if it is uniformly differentiable on $[a, b]$.
% \end{lem}
% \begin{proof}
% Uniform differentiability is defined as:
% \begin{equation*}
% \forall \epsilon > 0, \exists \delta > 0, \forall x, y \in [a, b]: |x - y| < \delta \implies \left|\frac{f(x) - f(y)}{x-y} - f'(y)\right| < \epsilon
% \end{equation*}
% $implies:$
% \end{proof}

\begin{lem}
If $f$ is of bounded variation on $[a, b]$, then $f$ is bounded on $[a, b]$.
\end{lem}
\begin{proof}
Consider some arbitrary $x \in [a, b]$. Then consider the simple partition $\Pi \coloneq \{a, x, b\}$. We see that the variation over that partition is:
\begin{equation*}
|f(x) - f(a)| + |f(b) - f(x)|
\end{equation*}
\begin{equation*}
\implies |f(x) - f(a)| \leqslant |f(x) - f(a)| + |f(b) - f(x)| \leqslant V(f, [a, b])
\end{equation*}
By the reverse triangle inequality, we obtain:
\begin{equation*}
\forall x \in [a, b]: |f(x)| \leqslant |f(a)| + V(f, [a, b]) 
\end{equation*}
\end{proof}

% \subsection{Total Variation and Stieltjes Integrals}
\subsection{Darboux Formulation}
\subsection{Henstock-Kurzweil Generalization}
\pagebreak
\section{Measure Theory*}
\pagebreak
\section{Arzela-Ascoli and Nets}
\end{document}
